\section{模型假设}
\subsection{模型基本假设}
\begin{enumerate}[label=(\arabic*), leftmargin=2.5em, itemsep=0.3em]
\item \textbf{假设应服务于建模。} 模型假设用于简化实际问题、明确研究范围，并为模型建立提供合理前提。

\item \textbf{假设必须有依据。 }可以忽略影响较小的因素，也可以将短期内变化不明显的量视为常数，但应符合实际情况。

\item \textbf{只写必要假设。} 不要罗列与模型无关、对求解没有影响的内容，也不要为了凑篇幅设置过多假设。

\item \textbf{假设应具体明确。} 避免使用“其他条件均不变”“不考虑特殊情况”等过于笼统的表述，应说明具体忽略了什么、保持什么不变。

\item \textbf{后文保持一致。} 模型建立、参数设置和结果分析必须遵循已提出的假设；若关键假设可能影响结果，应进行误差分析或灵敏度分析。
\end{enumerate}
\textbf{模型假设不是随意简化，而是在保证主要规律不变的前提下，忽略次要因素。}

\textbf{常见假设类型}

\ding{172} \textbf{忽略次要因素}
当某些因素对研究结果影响较小时，可以暂不考虑。
假设在研究时间范围内，天气变化对运输时间的影响可以忽略。

\ding{173} \textbf{将变量视为常数}
对于短期内变化较小或难以获取的数据，可在一定范围内视为常数。
假设研究期间产品单价和运输成本保持不变。

\ding{174} \textbf{理想化研究条件}
对实际过程进行合理理想化，以降低模型复杂度。
假设设备运行状态正常，不发生突发故障。

\ding{175} \textbf{排除小概率事件}
在题目未特别要求时，可不考虑极端情况或偶发事件。
假设研究期间不发生重大自然灾害和突发公共事件。

\textbf{写作提示}
\begin{enumerate}[label=(\arabic*), leftmargin=2.5em, itemsep=0.3em]
\item 每条假设尽量只表达一个意思，并采用编号分条书写。

\item 不要把题目已经明确给出的条件重复写成假设。

\item 避免使用明显不符合实际、只是为了方便计算的强假设。
\end{enumerate}


\textbf{写作原则：合理、必要、具体、可解释。}
